% TOP_PROBLEM: {"id": 3267, "title": "Arc-disjoint out-branching and in-branching", "category_id": 3, "category": {"id": 3, "name": "graph_theory", "display_name": "Graph Theory", "description": "Problems involving graphs, networks, and their properties.", "slug": "graph-theory", "order_index": 3, "created_at": "2026-07-31T15:26:25.671Z"}, "status": "open", "problem_number": "OPG-46495", "set_id": 12}
% FIRST_POSTED: 2026-10-01
% ENTRY_KIND: statement_only
% UnsolvedMath ID 3267; source catalogue code OPG-46495
% !TeX program = lualatex
% Definition and sources only; this file is not an LLM solution attempt.
\documentclass[11pt,a4paper]{article}
\usepackage[margin=25mm]{geometry}
\usepackage{fontspec}
\setmainfont{lmroman10-regular.otf}[BoldFont=lmroman10-bold.otf,
 ItalicFont=lmroman10-italic.otf,BoldItalicFont=lmroman10-bolditalic.otf]
\usepackage{amsmath,amssymb,mathtools}
\usepackage{unicode-math}
\setmathfont{latinmodern-math.otf}
\AtBeginDocument{\let\setminus\smallsetminus}
\usepackage{xurl,hyperref}
\hypersetup{colorlinks=true,urlcolor=blue}
\setcounter{secnumdepth}{0}
\setlength{\parindent}{0pt}
\setlength{\parskip}{5pt}
\setlength{\emergencystretch}{3em}
\begin{document}
\section*{Arc-disjoint out-branching and in-branching}\label{3267}
{\small UnsolvedMath ID 3267 · OPG-46495 · Graph Theory · OpenGarden}
\subsection{Definitions and mathematical statement}
A finite loopless digraph $D=(V,A)$ is $k$-arc-strong if deleting any set
of fewer than $k$ arcs leaves a strongly connected digraph. Strong
connectivity means a directed path exists from each vertex to each
other vertex. Equivalently, every nonempty proper subset $X\subset V$
has at least $k$ arcs leaving it. Arcs are counted individually.

An out-branching rooted at $u\in V$ is a spanning directed tree in which
$u$ has indegree zero, every other vertex has indegree one, and every
vertex is reachable from $u$. An in-branching rooted at $v\in V$ is the
reverse notion: every vertex can reach $v$, which has outdegree zero,
and every other vertex has outdegree one. Each branching uses exactly
$|V|-1$ arcs.

The conjecture asks whether some absolute integer $k\ge1$ has the
following property: for every finite $k$-arc-strong digraph $D$ and every
prescribed ordered pair $u,v\in V$, there exist an out-branching $B^+$
rooted at $u$ and an in-branching $B^-$ rooted at $v$ such that
\[
                              A(B^+)\cap A(B^-)=\varnothing.
\]
The roots may coincide. Both branchings necessarily use every vertex,
so the disjointness requirement concerns arcs only. The constant must
be independent of the order of $D$ and of the selected roots. Separate
existence of the two branchings does not settle the simultaneous
arc-disjoint requirement.
\subsection{Short English statement}
Does sufficiently high arc-connectivity guarantee arc-disjoint spanning outward and inward trees at any two prescribed roots?
\subsection{Sources}
\begin{itemize}
\item[S1] ULAM.AI and the UnsolvedMath Contributors, supplied problems.json, record 3267 (OPG-46495), OpenGarden collection. Catalogue identity and source transcription; CC BY 4.0.
\url{https://huggingface.co/datasets/ulamai/UnsolvedMath}.
\item[S2] Open Problem Garden, Arc-disjoint out-branching and in-branching. Original problem and discussion; the mathematical formulation below is expanded from the supplied source record.
\url{http://www.openproblemgarden.org/op/arc_disjoint_out_branching_and_in_branching}.
\end{itemize}
\end{document}
